\chapter{Literature Review}
\label{Literature Review}





{\color{red} \textbf{Instruction:} A literature review is a comprehensive examination and synthesis of existing scholarly works, including research articles, books, and other relevant sources, about a specific topic or research question. It plays a crucial role in academic writing by establishing the context of the research, identifying gaps and trends in the current knowledge landscape, and providing a theoretical framework for the study. Through a critical evaluation of the literature, it justifies the need for the proposed research, guides methodological choices, and helps avoid redundancy by building upon prior work. The literature review not only assesses the quality and credibility of existing research but also organizes information in a structured manner, categorizing studies based on themes or methodologies. Additionally, it serves as a roadmap for future research by highlighting unanswered questions and areas that warrant further investigation. Overall, a well-conducted literature review contributes to the depth and credibility of academic research, ensuring that the study is situated within the broader scholarly discourse.}











\section{Introduction}

{\color{blue}\textbf{Purpose:}
Set the scope and objectives of the literature review; tell the reader what topics, time span, and types of sources you will cover and why.} \\

{\color{red}\textbf{Writing Instructions (step-by-step):}
\begin{enumerate}
    \item Start with 2–3 short paragraphs that define the scope: subject area, timeframe (e.g., last 10 years), and types of literature (peer-reviewed journals, conference papers, standards, patents).
    
    \item Explicitly state the aims of the review (e.g., ``This review aims to synthesize recent empirical findings on X, evaluate conceptual models Y, and identify gaps relevant to Z'').
    
    \item Describe your search strategy briefly: databases searched (Scopus, IEEE Xplore, Web of Science), keywords, inclusion/exclusion criteria, and date range. Include a short PRISMA-style sentence or a tiny flow line if many sources.
    
    \item Outline how the review is organized (themes, theories, methods, or chronology). Use a short roadmap sentence: ``Section 2.2 reviews theoretical models; 2.3 surveys empirical studies; 2.4 critically compares approaches.''
\end{enumerate}}


\section{Theoretical Background / Existing Models}

{\color{blue}\textbf{Purpose:}
Present the core theories, conceptual models, and formal frameworks that underpin your research — provide the theoretical foundations you will use later.} \\

{\color{red}\textbf{Writing Instructions (step-by-step):}
\begin{enumerate}
    \item Identify and name the principal theories or models relevant to your problem (e.g., queuing theory, socio-technical models, diffusion of innovation, CNN architectures).
    
    \item For each theory/model: give a concise definition, its origin (key citations), formal structure (variables, equations, or components), and why it matters for your work. Use subheadings for each major theory.
    
    \item If applicable, include mathematical formulation or diagrams (labelled figures) to show model components and relationships. Explain notation clearly.
    
    \item Critically evaluate each model: strengths, assumptions, limitations, empirical support. Don't merely summarize — show critical appraisal.
    
    \item Conclude with how these theories inform your conceptual framework or methodological choices (bridge to Chapter 3).
\end{enumerate}}



\section{Related Studies}

{\color{blue}\textbf{Purpose:}
Review empirical and applied research closely related to your research question — demonstrate depth of reading and situate your work within the literature.} \\

{\color{red}\textbf{Writing Instructions (step-by-step):}
\begin{enumerate}
    \item Group studies by theme, method, or application rather than listing paper by paper. Use logical headings (e.g., ``Deep learning approaches,'' ``Rule-based systems,'' ``Field experiments in X'').
    
    \item For each group, synthesize key findings, datasets used, methodologies, evaluation metrics, and reported strengths and weaknesses. Use short comparative tables summarizing essential info (Author | Year | Method | Dataset | Key Result | Limitation).
    
    \item Highlight the most influential/closest works (2–5 anchor papers). Provide slightly deeper critique for these (design choices, parameter settings, reproducibility).
    
    \item Pay attention to methodological details: sample sizes, hyperparameters, baselines, evaluation protocols — note any inconsistencies across studies.
    
    \item Identify patterns — where the literature converges and diverges. Use transitional sentences to lead into gaps.
\end{enumerate}


\begin{itemize}
     \item Review existing literature related to your topic.
     \item Identify key theories, concepts, and research findings.
     \item Discuss the current state of knowledge in the field and gaps that your project/thesis aims to address.
\end{itemize} }


\textbf{ \color{blue} Example:}


\subsection*{Blockchain-Based IoT Access Control System: Towards Security, Lightweight, and Cross-Domain}
Sun et al. propose a blockchain-based access control system tailored for the Internet of Things (IoT) that ensures lightweight performance, cross-domain support, and enhanced security \cite{Sun2021}. The system employs smart contracts and a role-based permission structure to manage access in decentralized IoT environments. It introduces a lightweight consensus mechanism and uses elliptic curve cryptography to reduce overhead.




These limitation are a lack of performance validation under high node scalability. Energy consumption on resource-constrained IoT devices has not been fully evaluated. The paper does not address key revocation in access control.




\begin{longtable}{|p{3cm}|p{7cm}|p{5.0cm}|}
\caption{Literature review summary of blockchain-based certificate verification system} \\
\hline
\textbf{Author(s) (Year)} & \textbf{Contribution} & \textbf{Limitation} \\
\hline
\endfirsthead

\hline
\textbf{Author(s) (Year)} & \textbf{Contribution} & \textbf{Limitation} \\
\hline
\endhead

\hline
\endfoot

Rustemi et al. (2023) \cite{Rustemi2023} & Conducted a \textbf{systematic literature review} on blockchain-based academic certificate verification systems and classified smart contract usage patterns. & \textbf{Limited scope} to certificate verification; did not address \textbf{interoperability or legal compliance}. \\
\hline
KC et al. (2023) \cite{Aloqaily2020} & Developed a system using \textbf{QR code-based issuer validation} and \textbf{blockchain storage} to detect fake physical certificates. & Focused only on \textbf{PDF or printed certificates}; lacks support for \textbf{legacy institutional integration}. \\
\hline
Jadhav et al. (2024) \cite{Ankit2023} & Proposed a \textbf{blockchain-driven smart contract system} for secure issuance and verification of educational certificates. & \textbf{Limited testing} on large-scale deployment and \textbf{no discussion on user privacy or data anonymization}. \\
\hline
Rahman et al. (2023) \cite{Rahman2023} & Proposed an \textbf{AI-integrated blockchain model} for threat detection in educational networks. & \textbf{Too complex for typical academic deployment}; lacks certificate-specific validation logic. \\
\hline

\end{longtable}








\section{Comparative Analysis of Existing Approaches}

{\color{blue}\textbf{Purpose:}
Systematically compare principal approaches, highlighting performance, complexity, scalability, assumptions, and applicability to your problem domain.} \\

{\color{red}\textbf{Writing Instructions (step-by-step):}
\begin{enumerate}
    \item Choose comparison dimensions relevant to your research: accuracy/performance, data requirements, computational cost, generalizability, interpretability, robustness, ethical/privacy concerns.
    
    \item Create a clear, well-labelled comparison table (e.g., TABLE 2.1) with rows as approaches and columns as comparison criteria. Add footnotes for special cases.
    
    \item For each criterion, write a small paragraph interpreting the table: which approaches excel, which fail, and why. Use empirical evidence from literature (cite numbers where possible).
    
    \item Discuss trade-offs explicitly (e.g., approach A gives higher accuracy but requires large labeled data and high compute).
    
    \item Evaluate suitability for your research context and constraints — justify why certain approaches are promising or unsuitable.
\end{enumerate}}


\section{Summary of Research Gaps}

{\color{blue}\textbf{Purpose:}
Synthesize the review to state clearly what is missing in the literature and how your thesis will address these gaps.} \\

{\color{red}\textbf{Writing Instructions (step-by-step):}
\begin{enumerate}
    \item List 3–5 precise, researchable gaps derived from the comparative analysis (e.g., ``Lack of real-world longitudinal evaluation,'' ``Limited explainability in model X,'' ``No integration of privacy preserving methods in Y''). Phrase each gap as a concise sentence.
    
    \item For each gap, explain the consequences of that gap (why it matters) and cite supporting literature.
    
    \item Map gaps to potential research objectives — show how each gap suggests a specific aim or question you will pursue. A small mapping table is excellent: (Gap → Implication → Thesis Objective).
    
    \item End the section with a short paragraph explicitly stating how your study addresses the most critical gap(s) — this bridges to Chapter 3 (Conceptual Framework \& Methodology).
\end{enumerate}}



\textbf{ \color{blue} Example:}Traditional certificate verification methods mainly rely on centralized systems and manual checking, which often leads to document forgery and inefficiency \cite{Aloqaily2020, Rahman2023}. These processes take a lot of time and cause delays in important activities such as job recruitment, student admissions, and other tasks that require quick verification. Centralized databases are also at risk of tampering, forgery, and data manipulation, which reduces transparency and trust \cite{Rahman2023}. In Bangladesh, the verification of myGov.bd certificates are done online, but it still depends on a centralized database without the use of blockchain technology, which limits transparency and resistance to tampering \cite{Rustemi2023} . Globally, there is no universal, secure, and fast method for certificate verification \cite{Ankit2023}. Therefore, there is a clear need for a decentralized, tamper-proof solution to ensure the authenticity, integrity, and quick verification of certificates  \cite{Sun2021}.